\section{The problem and the precise improvement}
\label{sec:scope}

The coefficients
\begin{equation}\label{eq:definition-a}
 a_j=(j+1)\log(j+1)-j\log j-1
     =\int_j^{j+1}\log x\,dx\qquad(j\ge1)
\end{equation}
are positive and grow like $\log j$. Set
\begin{equation}\label{eq:definition-g}
 A(z)=\sum_{j\ge1}a_jz^j,\qquad
 g_0=1,\qquad g_n=\sum_{j=1}^n a_jg_{n-j},\qquad
 G(z)=\frac1{1-A(z)}.
\end{equation}
The radius of convergence of $A$ is one. On $[0,1)$ it increases
strictly from zero to infinity. Thus there is a unique $\rho\in(0,1)$
with $A(\rho)=1$. Throughout the paper,
\begin{equation}\label{eq:normalization}
 \gamma=\frac1{\rho A'(\rho)},\qquad
 w_0=\rho^{-1},\qquad
 E_n=\gamma\rho^{-n}-g_n\quad(n\ge1).
\end{equation}
Euler's constant is always $\gE$, to distinguish it from the residue
amplitude $\gamma$. All logarithms are natural. Numerical values, given
only for orientation, are
\[
 \rho=0.542598586098471021959384595779\ldots,
 \qquad \gamma=0.323457534191380676574821548566\ldots.
\]

\subsection{Published baseline and repository connection}

Ford's study of totients uses this recurrence in its simplex-volume
estimates. Lemma~3.7 of the revised paper proves
$|g_n-\gamma\rho^{-n}|\le5$; the following remark records that the
errors for $n\ge1$ are negative, increasing, and summable
\cite[Lemma~3.7 and Remark~1]{Ford}. Ford attributes the simplified
recurrence argument to a solution with Lau \cite{FordLau}.
The present paper starts from those established facts as a comparison;
it proves the new estimates directly.

The same recurrence is defined in the inspected \texttt{openai/math}
totient manuscript, where a continuous counting scale contains
\begin{equation}\label{eq:counting-scale}
 \mathcal G_m(B)=\frac{B^m}{m!\prod_{j=1}^m g_j}.
\end{equation}
We use that definition and its motivation, without using the manuscript's
claims about the arithmetic counting function \cite{OpenAIMath}.
The study also fits the asymptotic-expansion and inversion program in
\texttt{ProveIt} \cite{ProveIt}. Section~\ref{sec:products} gives a
direct, proved consequence for \eqref{eq:counting-scale}.

The priority review inspected Ford's full revised paper and the relevant
repository source. The Ford--Lau item was identified bibliographically,
but its full text was not accessible during preparation. Targeted
searches did not locate the explicit remainder expansion proved here.
This supports a precise claim of a refinement of the inspected published
estimate; it does not establish worldwide first-publication priority.
Spectral measures, Hausdorff moments, renewal generating functions, and
polylogarithm expansions are classical tools.

\subsection{Main conclusions at a glance}

\begin{theorem}[Sharp remainder and all logarithmic orders]
\label{thm:main-log}
For every fixed integer $K\ge0$, as $n\to\infty$,
\begin{equation}\label{eq:main-log}
 E_n=\frac1{n^2}\left\{
    \sum_{k=0}^K\frac{c_k}{(\log n)^{k+2}}
      +O_K\bigl((\log n)^{-K-3}\bigr)\right\},
\end{equation}
where
\begin{equation}\label{eq:main-coeff}
 c_k=(k+1)![u^k]\,
           \frac{(1+u)e^{\gE u}}{\Gamma(1-u)}.
\end{equation}
In particular,
\begin{align}\label{eq:first-coefficients}
 c_0&=1,&c_1&=2,&c_2&=-\frac{\pi^2}{2},\nonumber\\
 c_3&=-2\pi^2-8\zeta(3),&
 c_4&=\frac{\pi^4}{12}-40\zeta(3).
\end{align}
Moreover, $E_n>0$ for every $n\ge1$, and for every $n\ge4096$,
\begin{equation}\label{eq:main-explicit}
 \frac1{200n(n+1)\log^2n}\le E_n
           \le\frac{10}{n(n+1)\log^2n}.
\end{equation}
\end{theorem}

Thus the formerly bounded absolute error actually tends to zero, with
an exact leading constant and an optimal power of $n$ and $\log n$.
After just the first correction,
\[
 g_n=\gamma\rho^{-n}-\frac1{n^2\log^2n}
                   +O\!\left(\frac1{n^2\log^3n}\right).
\]
The new $O$-term is itself sharp: its leading coefficient is $-2$.
The more conservative constants in \eqref{eq:main-explicit} provide
a statement valid from a specified finite index, independently of
any numerical calculation of $\rho$.

\begin{theorem}[Exact moment structure]
\label{thm:main-moment}
There is an explicit positive continuous function $\ell$ on $(0,1)$
such that
\begin{equation}\label{eq:main-moment}
 E_n=\int_0^1 t^{n-1}\ell(t)\,dt\qquad(n\ge1).
\end{equation}
Consequently $(-1)^r\Delta^r E_n>0$ for all $n\ge1$, $r\ge0$,
and every matrix $(E_{a+i+j})_{0\le i,j\le d}$ with $a\ge1$ is
positive definite. The exact sum rules are
\begin{equation}\label{eq:main-sums}
 \sum_{n\ge1}E_n=1-\frac\gamma{1-\rho},\qquad
 \sum_{n\ge1}nE_n=\frac{\gamma\rho}{(1-\rho)^2}.
\end{equation}
\end{theorem}

The restriction $n\ge1$ matters: $g_0-\gamma=1-\gamma>0$ has the
opposite sign. The proof explicitly checks both endpoints of the
spectral measure and excludes an unaccounted singular component.

\subsection{A second, finer expansion}

A fixed truncation of \eqref{eq:main-log} has an error much larger
than $n^{-3}$. To resolve powers of $1/n$ as well, define the ordinary
convergent integrals
\begin{equation}\label{eq:def-Fa}
 F_a(L)=\int_0^\infty
 \frac{y^{a-1}e^{-y}}{(L-\log y-\gE)^2+\pi^2}\,dy
 \qquad(a>0,\ L\in\mathbb R),
\end{equation}
and put
\begin{equation}\label{eq:p2p3}
 p_2=\zeta'(-1)-\frac1{12},\qquad p_3=\zeta'(-2).
\end{equation}
Primes on $F_a$ mean differentiation in $L$.

\begin{theorem}[Expansion retaining all logarithmic dependence]
\label{thm:main-sectors}
Writing $L=\log n$, one has
\begin{align}\label{eq:main-sectors}
 E_n={}&\frac{F_2(L)}{n^2}+\frac{F_3(L)}{2n^3}
   +\frac1{n^4}\left\{\frac{F_4(L)}{12}-p_2F_4'(L)\right\}\nonumber\\
 &-\frac1{n^5}\left(\frac{p_2}{2}+p_3\right)F_5'(L)
   +O\!\left(\frac1{n^6\log^2n}\right).
\end{align}
More generally, an explicit recursion generates the coefficient of
every power $n^{-r-2}$, with remainder
$O_R(n^{-R-3}\log^{-2}n)$ after powers through $n^{-R-2}$.
\end{theorem}

In \eqref{eq:main-sectors}, the $F_a$ are retained as integrals.
Replacing them by a fixed number of inverse-logarithmic terms would
destroy the displayed precision. This distinction is essential when
using the result as an asymptotic expansion in more than one scale.

\subsection{What these theorems do and do not settle}

The main theorem improves a precise analytic component of the totient
method. It also supplies exact inequalities for the entire recurrence,
not merely its tail. The product theorem below determines the first
relative correction to \eqref{eq:counting-scale}; the inverse theorem
quantifies convergence of the associated renewal process.

These conclusions do not, by themselves, improve the error term in an
asymptotic formula for the number of distinct totients. The arithmetic
steps that connect a volume calculation to a count have separate errors.
The present argument neither verifies nor requires the larger claims
in the repository manuscript. This boundary keeps the strengthened
estimate usable as an independently checkable lemma.
